% Preparación autónoma de VII, 2026-09-30. Fragmento sin inputs adicionales.
% RESULTADOS_RECUPERADOS, pruebas reunidas sin modificación de los originales:
% /Users/ruben/Documents/excelencia academica/output/CUMPLIMIENTO_EDITORIAL_20260928/fuentes/VIII_ES/manuscrito/29_retorno_y_memoria_traslacional.tex
% /Users/ruben/Documents/excelencia academica/output/CUMPLIMIENTO_EDITORIAL_20260928/fuentes/VIII_ES/manuscrito/30_pantalla_y_respuesta.tex
% /Users/ruben/Documents/excelencia academica/output/CUMPLIMIENTO_EDITORIAL_20260928/fuentes/VIII_ES/manuscrito/30b_variacion_energia_memoria.tex
% /Users/ruben/Documents/excelencia academica/output/CUMPLIMIENTO_EDITORIAL_20260928/fuentes/VIII_ES/manuscrito/30c_composicion_corriente_conexion.tex
% /Users/ruben/Documents/excelencia academica/output/CUMPLIMIENTO_EDITORIAL_20260928/fuentes/VIII_ES/manuscrito/30d_realizacion_geometrica.tex
% /Users/ruben/Documents/excelencia academica/output/CUMPLIMIENTO_EDITORIAL_20260928/fuentes/VIII_ES/manuscrito/31_corriente_y_respuesta_cartan.tex
% /Users/ruben/Documents/excelencia academica/output/CUMPLIMIENTO_EDITORIAL_20260928/fuentes/VIII_ES/manuscrito/33_corriente_espinorial_y_densidad.tex
% Dependencias internas ya presentes: núcleo común, rutas, lectores y unidades.
% Las inversas de 31 están demostradas en 10, hmtcuatro:gea:torsion.
% Legendre, fuente completa y restricciones permanecen en 20.
% Se incluye de 33 sólo lo utilizado por 10+20, no sus estados cosmológicos.
\section{Geometric and material background of the joint action}
\label{hmtcuatro:ant:seccion}

The seeds and operations of APP, the orientation and regime of TRIT, and the transport of TPK produce the enriched routes on which the readers of the common core act. Sheet, orientation, carry, boundary and memory are jointly preserved; the phase return does not reset that state. The constructions of the continuum belong to the same object and are not turned here into independent sectors. The realization that follows transports this already constructed structure to geometry and matter. Its units and couplings are earlier HMT outputs: no numerical comparison acts as a generator.

To make the action and its reduction readable within this manuscript, we bring together the screen operations, the effective variation of transport, and the spinorial current they use. The Cartan--Holst elimination is proved in \ref{hmtcuatro:gea:torsion}; the total action and its canonical reduction, in \ref{hmtcuatro:accion-retroaccion-restricciones}. The formulas in this section specify the data and calculations entering both.

\subsection{Incidence, screen and soldering}
\label{hmtcuatro:ant:pantalla}

The cubic representation of a route boundary has the six faces \(F=\{(i,\epsilon):1\leq i\leq3,\ \epsilon=\pm1\}\). On \(H_F=\mathbb R^F\), opposition is \(Re_{i,\epsilon}=e_{i,-\epsilon}\). With
\[
 u=\sum_{i,\epsilon}e_{i,\epsilon},\qquad
 P_u=\frac{uu^{\mathsf T}}6,\qquad
 P_-=\frac{I-R}2,\qquad
 P_0=\frac{I+R}2-P_u ,
\]
the three projectors are orthogonal, sum to \(I\) and have ranks \(1,3,2\). The incidence between faces and vertices \(\nu\in\{-1,1\}^3\) is \(M_{(i,\epsilon),\nu}=\mathbf1_{\{\nu_i=\epsilon\}}\).

\begin{proposition}[Screen of oriented incidence]
\label{hmtcuatro:ant:prop-pantalla}
We have
\[
 MM^{\mathsf T}=12P_u+4P_-,\qquad
 \ker M^{\mathsf T}=P_0H_F .
\]
The quotient \(Q=H_F/P_0H_F\), identified with \((P_u+P_-)H_F\), is four-dimensional. The form \(B=P_--P_u\) on \(Q\) has signature \((-+++)\).
\end{proposition}
\begin{proof}
Each face contains four vertices; two opposite faces share none, and two faces on different axes share two. Therefore,
\[
 MM^{\mathsf T}=4I+2(uu^{\mathsf T}-I-R)
              =12P_u+4P_- .
\]
The identity \(\|M^{\mathsf T}x\|^2=\langle x,MM^{\mathsf T}x\rangle\) determines the kernel. In the orthonormal basis
\[
 t=u/\sqrt6,\qquad d_i=(e_{i,+}-e_{i,-})/\sqrt2
\]
of \(Q\), the matrix of \(B\) is \(\operatorname{diag}(-1,1,1,1)\).
\end{proof}

Let \(W_\square=\sqrt6P_u+\sqrt2P_-\). Then
\[
 W_\square e_{i,\epsilon}=t+\epsilon d_i=:n_{i,\epsilon},
 \qquad B(n_{i,\epsilon},n_{i,\epsilon})=0,\qquad
 MM^{\mathsf T}|_Q=2W_\square^*W_\square|_Q .
\]
These equalities are obtained by applying the projectors to the face basis. The proper rotations of the cube preserve the projectors, the time orientation \(t\) and the form \(B\).

The enriched route provides each edge \(a\) with its face, its sign \(\sigma_m(a)\) and the transport \(U_m(a):Q_{s(a)}\to Q_{t(a)}\). Its cellular soldering is
\begin{equation}
 \beta_m(a)=\ell_m\sigma_m(a)n_{\lambda(\underline a)}^{s(a)},
 \quad \ell_m=\ell_0\,9^{-m},\quad
 \beta_m(\bar a)=-U_m(a)\beta_m(a).
 \label{hmtcuatro:ant:soldadura}
\end{equation}
The inversion \(U_m(\bar a)=U_m(a)^{-1}\) implies that inverting twice recovers the soldering. The ordered carry record is also preserved: it is not replaced by the arrival vector.

\begin{proposition}[Naturality of soldering]
\label{hmtcuatro:ant:refinamiento-soldadura} Let \(a=a_1\cdots a_9\) be a compatible refinement: the child incidences transported to the origin of the predecessor preserve face and sign, and \(\ell_{m+1}=\ell_m/9\). If \(\mathcal U_j:Q_{m,s(a)}\to Q_{m+1,s(a_j)}\) includes the identification by refinement and the transport to \(s(a_j)\), then
\[
 \beta_m(a)=\sum_{j=1}^9\mathcal U_j^{-1}\beta_{m+1}(a_j).
\]
\end{proposition}
\begin{proof}
Each summand in the source fibre is \((\ell_m/9)\sigma_m(a)n_{\lambda(\underline a)}\). Their sum gives \eqref{hmtcuatro:ant:soldadura}. Reversing the route reverses the order of the transports and each transport; the soldering inversion formula gives the global sign. Concatenation preserves the memory of the nine incidences in the same order. The equality iterates by composition of compatible refinements.
\end{proof}

\subsection{Transport representation and geometric realization}
\label{hmtcuatro:ant:geometria}

The record of a route \(\gamma\) counts inversions and effective changes of sheet through the cocycle \(c(\gamma)=(c_g(\gamma),c_s(\gamma))\in\mathbb Z^2\). The count is cumulative on forward routes; its extension to the groupoid is oriented and assigns opposite increments to the inverse route. For composable routes, \(c(\gamma_2\gamma_1)=c(\gamma_2)+c(\gamma_1)\), and \(c(\bar\gamma)=-c(\gamma)\). Let \(R_4\) be an element of order four in \(\operatorname{Spin}^+(3,1)\) whose vector representation \(\Lambda(R_4)\) fixes the unit time vector \(u_0\). The received affine representation is
\[
 \rho_{\rm C}^{\rm disc}(\gamma)
 =\bigl(R_4^{c_g(\gamma)},\ell_0c_s(\gamma)u_0\bigr).
\]
With the product \((R,a)(S,b)=(RS,a+\Lambda(R)b)\), the fixing of \(u_0\) and the additivity of \(c\) directly prove that \(\rho_{\rm C}^{\rm disc}\) preserves identity, concatenation and inversion. When the rotational component returns to the identity, the memory translation remains determined by \(c_s(\gamma)\), which is not reset.

Normalization
\[
 N_{\ell_0}(R,a)=(\Lambda(R),a/\ell_0),\qquad
 \rho_{\rm C}=N_{\ell_0}\circ\rho_{\rm C}^{\rm disc}
\]
is also a homomorphism, since \((a+\Lambda(R)b)/\ell_0=a/\ell_0+\Lambda(R)(b/\ell_0)\). The lift \(R_4^{c_g(\gamma)}\) is retained when spinors act; it cannot be recovered from \(\Lambda(R_4^{c_g(\gamma)})\) alone, whose representation has kernel \(\{1,-1\}\).

In a screen realization \(J_{{\rm scr},m}\) that transports the form and the soldering, \(e_m=J_{{\rm scr},m}\beta_m\). The smooth realization considered consists of an oriented four-dimensional region \(\mathcal U\), a nondegenerate cotetrad \(e\), a Lorentz connection \(\omega\) and a route realization \(\mathfrak R_{\rm C}\) with the compatibility condition
\begin{equation}
 \operatorname{Hol}_{\mathcal A_{\ell_0}}
          (\mathfrak R_{\rm C}(\gamma))
 =\rho_{\rm C}(\operatorname{Hol}_{\rm TPK}(\gamma)),\qquad
 \mathcal A_{\ell_0}=
 \begin{pmatrix}\omega&\ell_0^{-1}e\\0&0\end{pmatrix}.
\label{hmtcuatro:ant:compatibilidad-holonomia}
\end{equation}
The argument of \(\rho_{\rm C}\) preserves the cocycle of the enriched holonomy; it is not replaced by its reduced phase alone. The scale \(\ell_0>0\) is constant in the coordinatization. These are the data and the condition of the smooth realization used; preserving path composition neither identifies distinct paths with the same endpoints nor removes their memory. The metric is \(g=\eta_{IJ}e^I\otimes e^J\), \(\eta=\operatorname{diag}(-1,1,1,1)\).

\begin{proposition}[Curvature, covariance and Bianchi]
\label{hmtcuatro:ant:curvatura} With \(F_\omega=\mathrm d\omega+\omega\wedge\omega\) and \(T=\mathrm de+\omega\wedge e\),
\[
 \mathcal F_{\ell_0}=
 \begin{pmatrix}F_\omega&\ell_0^{-1}T\\0&0\end{pmatrix},
 \qquad D_\omega T=F_\omega\wedge e,\quad D_\omega F_\omega=0 .
\]
A change of frame \(g_L:\mathcal U\to SO^+(3,1)\) acts by \(e'=g_Le\), \(\omega'=g_L\omega g_L^{-1}-\mathrm dg_L\,g_L^{-1}\); then \(T'=g_LT\), \(F_{\omega'}=g_LF_\omega g_L^{-1}\).
\end{proposition}
\begin{proof}
Multiplying matrices of forms gives the diagonal block \(\mathrm d\omega+\omega\wedge\omega\) and the translational block \(\ell_0^{-1}(\mathrm de+\omega\wedge e)\). On substituting the transformed fields, the terms involving \(\mathrm dg_L\) cancel. Expanding \(D_\omega T\) leaves \((\mathrm d\omega+\omega\wedge\omega)\wedge e\). In \(D_\omega F_\omega\), the derivatives of \(\omega\) and the cubic products cancel in pairs.
\end{proof}

The tritic sheet \(\tau\in\{+1,0,-1\}\) preserves its regime in the Cartan realization with generators \(J_{ab},P_a^{(\tau)}\), the Lorentz vector action and \([P_a^{(\tau)},P_b^{(\tau)}]=-\tau J_{ab}\). For \(\mathcal A_\tau=\frac12\omega^{ab}J_{ab}+\ell^{-1}e^aP_a^{(\tau)}\),
\[
 \mathcal F_\tau=
 \tfrac12(F_\omega^{ab}-\tau\ell^{-2}e^a\wedge e^b)J_{ab}
 +\ell^{-1}T^aP_a^{(\tau)} .
\]
Indeed, the Lorentz bracket gives \(F_\omega\), the mixed bracket gives \(D_\omega e\) and the translational bracket gives the term \(-\tau e\wedge e/(2\ell^2)\). The condition \eqref{hmtcuatro:ant:compatibilidad-holonomia} corresponds to the central sheet; the extreme sheets use their own representation \(\rho_{{\rm C},\tau}\) and holonomy condition.

\subsection{Transport differential and material current}
\label{hmtcuatro:ant:corriente-transporte}

In a finite cellular realization, let \(U_a:\mathcal H_{s(a)}\to\mathcal H_{t(a)}\) be an invertible transport between Hermitian fibres. It is not assumed that every Lorentz transport is unitary for the energy inner product. One orientation per edge is chosen. For
\[
 (Df)_a=f_{t(a)}-U_af_{s(a)},\quad r=Df,\quad
 q=\mathsf W r,\quad v_a=U_af_{s(a)},\quad
 E=\langle r,\mathsf W r\rangle ,
\]
the weight \(\mathsf W=\mathsf W^*>0\) may couple edges. This weight is not \(W_\square\). The variations correspond to the connection representation and the chosen boundary problem.

\begin{proposition}[Full variation]
\label{hmtcuatro:ant:variacion-completa} For a differentiable curve of data, with \(A_a=\dot U_aU_a^{-1}\), we have
\begin{equation}
 \dot E=2\operatorname{Re}\sum_a
 \langle q_a,\dot f_{t(a)}-U_a\dot f_{s(a)}-A_av_a\rangle
 +\langle r,\dot{\mathsf W}r\rangle .
 \label{hmtcuatro:ant:variacion}
\end{equation}
With the field and weight fixed, the full transport gradient is \(G_a=-2q_av_a^*\):
\[
 \delta_U E=\sum_a\operatorname{Re}\operatorname{Tr}(G_a^*A_a).
\]
In the unitary subclass \(A_a^*=-A_a\), its anti-Hermitian part \(J_a=v_aq_a^*-q_av_a^*\) suffices.
\end{proposition}
\begin{proof}
The product rule gives \(\dot E=2\operatorname{Re}\langle\mathsf W r,\dot r\rangle+
\langle r,\dot{\mathsf W}r\rangle\). Differentiating each residual yields \eqref{hmtcuatro:ant:variacion}. Since \(\operatorname{Tr}(v_aq_a^*A_a)=q_a^*A_av_a\), we obtain \(G_a\). The Hermitian part of \(G_a\) is orthogonal, with respect to \(\operatorname{Re}\operatorname{Tr}(X^*Y)\), to the anti-Hermitian variations. This reduction applies only in that subclass.
\end{proof}

For a route with \(U_\gamma=U_N\cdots U_1\), let \(B_k=U_N\cdots U_{k+1}\). The product rule proves
\begin{equation}
 \dot U_\gamma U_\gamma^{-1}
 =\sum_k B_kA_kB_k^{-1}.
 \label{hmtcuatro:ant:variacion-ruta}
\end{equation}
Consequently, a gradient \(G_\gamma\) is transported to the edge \(k\) as \(G_k=B_k^*G_\gamma(B_k^{-1})^*\). Cyclicity of the trace proves this formula even if \(B_k\) is not unitary.

If the real coordinates of connection variation are \(\theta_b\), its effective differential is written as \(A_a(\theta)=\sum_b B_{ab}\theta_b\). For \(S_{\rm mem}=\mathfrak a E\), with the field, weight and action section \(\mathfrak a\) fixed,
\begin{equation}
 \delta S_{\rm mem}=-\tfrac12\sum_b\theta_b s_b,\qquad
 s_b=4\mathfrak a\,\operatorname{Re}
          \sum_a\langle q_a,B_{ab}v_a\rangle .
 \label{hmtcuatro:ant:corriente-celular}
\end{equation}
Substituting \(A_a(\theta)\) into the variation proves the equality and fixes the factor \(4\). If the oriented cellular pairing is an invertible matrix \(\mathsf M\), defined by \(\delta S_{\rm mem}=-\frac12\theta^{\mathsf T}\mathsf M\sigma\), the current is \(\sigma=\mathsf M^{-1}s\). This matrix preserves volumes and orientation; it cannot be replaced by scalar division without fixing a basis in which that is appropriate. If \(f,\mathsf W,\mathfrak a\) also vary, all the terms of \eqref{hmtcuatro:ant:variacion} and \(E\delta\mathfrak a\) are retained.

For a weight per edge, inversion preserves the full response by jointly transforming
\[
 U_{\bar a}=U_a^{-1},\qquad
 r_{\bar a}=-U_a^{-1}r_a,\qquad
 \mathsf W_{\bar a}=U_a^*\mathsf W_aU_a .
\]
Substitution gives \(r_{\bar a}^*\mathsf W_{\bar a}r_{\bar a}=r_a^*\mathsf W_ar_a\). Its derivative necessarily retains \(\dot{\mathsf W}_{\bar a}
=U_a^*(A_a^*\mathsf W_a+\dot{\mathsf W}_a+\mathsf W_aA_a)U_a\): removing this term would alter the current. If the weight couples several edges, the same congruence is applied to the full matrix with block-diagonal transport; its cross-blocks are also retained.

\begin{proposition}[Variational compatibility under refinement]
\label{hmtcuatro:ant:refinamiento-corriente} For fixed maps \(J_m,K_m\), suppose that, throughout the variation considered,
\[
 D_{m+1}J_m=K_mD_m,\qquad
 K_m^*\mathsf W_{m+1}K_m=\mathsf W_m .
\]
Then \(E_{m+1}(J_mf)=E_m(f)\) and their differentials coincide. If the actions also coincide and \(P_m\) transports the connection variations to the refinement, their currents satisfy
\[
 s_m=P_m^{\mathsf T}s_{m+1},\qquad
 \mathsf M_m\sigma_m=P_m^{\mathsf T}\mathsf M_{m+1}\sigma_{m+1}.
\]
\end{proposition}
\begin{proof}
Substituting \(D_{m+1}J_m=K_mD_m\) into the energy and then the weight congruence proves the first equality along the entire curve. Differentiating it and using \(-\frac12\theta^{\mathsf T}s_m
=-\frac12(P_m\theta)^{\mathsf T}s_{m+1}\) proves the others. If \(J_m\) varies, its differential contributes \(2\operatorname{Re}\langle D_{m+1}J_mf,
\mathsf W_{m+1}D_{m+1}\dot J_m f\rangle\); it is not discarded.
\end{proof}

The memory current thus obtained comes from its action. The following spinorial current comes from the affine action of the fields. They combine when they are terms of a single action; they are not identified merely because they share transport. In particular, eliminating a contorsion on which the prepared field, the incidence or its minimum also depend requires retaining those differentials. The linear Cartan formula used below corresponds to the expressly stated affine material problem.

\subsection{Spinorial representation and affine current}
\label{hmtcuatro:ant:espinores}

The spin structure of the geometric realization is fixed. The inherited matrices, after complexifying the real module, are
\begin{equation}
 \begin{gathered}
 J=\begin{pmatrix}0&-1\\1&0\end{pmatrix},\quad
 R=\begin{pmatrix}0&1\\1&0\end{pmatrix},\quad
 Z=\begin{pmatrix}1&0\\0&-1\end{pmatrix},\\
 g^0=J\otimes I_2,\quad g^1=R\otimes I_2,\quad
 g^2=Z\otimes R,\quad g^3=Z\otimes Z .
 \end{gathered}
 \label{hmtcuatro:ant:clifford}
\end{equation}
Their squares are \(-I_4,I_4,I_4,I_4\), and distinct matrices anticommute. It follows that \(\{g^I,g^J\}=2\eta^{IJ}I_4\). We define
\[
 g_I=\eta_{IJ}g^J,\quad
 \mathbb B_{IJ}=\tfrac14[g_I,g_J],\quad
 A_D=-ig^0,\quad \bar\psi=\psi^\dagger A_D,\quad
 \Omega=\tfrac12\omega^{IJ}\mathbb B_{IJ}.
\]
Expanding the three-matrix commutator gives \([\mathbb B_{IJ},g^K]=\delta_J^Kg_I-\delta_I^Kg_J\). Thus \(\Omega\) represents the same metric connection, with \(D\psi=\mathrm d\psi+\Omega\psi\) and \(D\bar\psi=\mathrm d\bar\psi-\bar\psi\Omega\). The matrix \(A_D\) is Hermitian and \(A_Dg^I\) is anti-Hermitian; these fix the adjoint and the action factor in signature \((-+++)\).

Let \(\eta_I=\iota_{E_I}\operatorname{vol}_e\), so that \(e^J\wedge\eta_I=\delta_I^J\operatorname{vol}_e\). The affine action is
\begin{equation}
 S_D=\int_M\left\{
 \frac{\hbar}{2}
 [\bar\psi g^ID\psi-(D\bar\psi)g^I\psi]\wedge\eta_I
 -mc\,\bar\psi\psi\operatorname{vol}_e\right\}.
 \label{hmtcuatro:ant:accion-dirac}
\end{equation}
The mass and the positive sections \(c,\hbar\) are the received readings, and \([\psi]=L^{-3/2}\). The alternative notation \(\Gamma^I=-ig^I\), with principal term \(i\hbar\Gamma^ID_I\), switches to the opposite Clifford signature; the two conventions are not mixed.

\begin{proposition}[Normalized spinorial current]
\label{hmtcuatro:ant:corriente-spin} With \(B^{IJK}=\bar\psi g^{[I}g^Jg^{K]}\psi\), where antisymmetrization is normalized, the convention \(\delta_\omega S_D=-\frac12\int\delta\omega^{JK}\wedge\sigma_{JK}\) gives
\begin{equation}
 \sigma_{JK}
 =-\frac{\hbar}{2}\bar\psi\{g^I,\mathbb B_{JK}\}\psi\,\eta_I
 =-\frac{\hbar}{2}B^I{}_{JK}\eta_I .
 \label{hmtcuatro:ant:sigma-spin}
\end{equation}
The current is independent of the contorsion when \(e,\psi\) is fixed.
\end{proposition}
\begin{proof}
We have \(\delta\Omega=\frac12\delta\omega^{JK}\mathbb B_{JK}\). The two kinetic terms of \eqref{hmtcuatro:ant:accion-dirac} give
\[
 \delta_\omega S_D=\frac{\hbar}{4}\int
 \delta\omega^{JK}\wedge\eta_I\,
          \bar\psi\{g^I,\mathbb B_{JK}\}\psi .
\]
Comparison with the variational convention determines the sign and factor. Expanding the anticommutator using Clifford relations cancels the degree-one contractions and leaves \(\{g^I,\mathbb B_{JK}\}=g^{[I}g_Jg_{K]}\). The action is linear in \(\Omega\), which proves independence.
\end{proof}

If internal connections are included, they act on the material factor of the spinorial module, whereas \(\mathbb B_{IJ}\) acts on the spin factor. Thus, \([\,\mathbb B_{IJ}\otimes I,I\otimes\rho(A)\,]=0\). Varying \(\omega\) produces the same formula, contracted over the internal components and summed over the multiplets present. This preserves their chiral, colour and family representations; it adds no multiplets to those already constructed.

\subsection{Torsional contraction and cross-terms between species}
\label{hmtcuatro:ant:contraccion}

The real inversion of Cartan--Holst and of the map \(T^I\mapsto T^I\wedge e^J-e^I\wedge T^J\), proved in \ref{hmtcuatro:gea:torsion}, unambiguously fixes the contorsion \(\mathfrak k_*\) for the affine current \eqref{hmtcuatro:ant:sigma-spin}. These formulas retain \(\gamma\in\mathbb R\setminus\{0\}\), a nondegenerate cotetrad, and the sections \(G,c,\hbar,\gamma\) fixed during variation. Elimination uses the linear interaction only once and gives \(-\frac14\mathfrak k_*^{IJ}\wedge\sigma_{IJ}\).

\begin{theorem}[Coefficient of the spinorial contraction]
\label{hmtcuatro:ant:coeficiente} Let \(\kappa=8\pi G/c^4\). The effective interaction is
\[
 \mathcal L_{\rm spin}^{\rm ef}
 =C_\gamma q(B)\operatorname{vol}_e,\qquad
 C_\gamma=\frac{3\kappa c\hbar^2}{16}
                    \frac{\gamma^2}{1+\gamma^2},
\]
where
\[
 q(B)=\frac1{3!}B^{IJK}B_{IJK}
 =-(B^{012})^2-(B^{013})^2-(B^{023})^2+(B^{123})^2 .
\]
\end{theorem}
\begin{proof}
We first calculate
\[
 \mathcal Y^{IJ}
 =\kappa c\frac{\gamma^2}{1+\gamma^2}
       (-\star+\gamma^{-1}I)\sigma^{IJ},\quad
 \mathcal B^I=\sum_J\iota_{E_J}\mathcal Y^{IJ},\quad
 T^I=\mathcal B^I-\tfrac14e^I\wedge\sum_L\iota_{E_L}\mathcal B^L,
\]
and then \(\mathfrak k_{IJK}=(T_{KIJ}-T_{IJK}-T_{JKI})/2\). These are linear compositions of the already proved inverses. To make all the contraction coefficients explicit, the trivector is ordered as \((012,013,023,123)\) and the factor \(\kappa c\hbar^2\gamma^2/(1+\gamma^2)\) is extracted. Substituting \(\sigma_{JK}/\hbar=-\frac12B^I{}_{JK}\eta_I\) into these compositions and into \(-\frac14\mathfrak k^{IJ}\wedge\sigma_{IJ}\) gives
\[
 \begin{array}{c|rrrr|c}
  \text{operator on }\sigma&012&013&023&123&
           \text{six polarization coefficients}\\ \hline
  -\star&-3/16&-3/16&-3/16&3/16&0\\
  I&0&0&0&0&0
 \end{array}
\]
Each entry uses \(e^I\wedge\eta_J=\delta_J^I\operatorname{vol}_e\) and the contractions written above. The four diagonal entries and the six polarizations enumerate the full quadratic form. The part \(\gamma^{-1}I\) vanishes and the remainder is \((3/16)q(B)\), which proves the formula.
\end{proof}

For several species, variation is performed on their common action: \(\sigma=\sum_s\sigma_s\). Since the response \(\mathfrak k_*=\mathsf C_e(\sigma)\) is linear,
\begin{equation}
 -\tfrac14\mathsf C_e\!\left(\sum_s\sigma_s\right)
       \wedge\sum_t\sigma_t
 =-\tfrac14\sum_{s,t}\mathsf C_e(\sigma_s)\wedge\sigma_t .
 \label{hmtcuatro:ant:cruces-especies}
\end{equation}
This expansion proves that the terms \(s\ne t\) remain. Calculating separate responses and retaining only \(s=t\) would change the total action.

\subsection{Normal ordering of the total current}
\label{hmtcuatro:ant:car}

The finite fermionic realization of a cell uses \(\mathcal F=\bigoplus_n\Lambda^n V\), with the relations \(\{a_r,a_s^\dagger\}=\delta_{rs}\), \(\{a_r,a_s\}=0\) and the occupation vacuum as the normal-ordering reference. We write \(d\Gamma(M)=\sum_{r,s}a_r^\dagger M_{rs}a_s\); this \(\Gamma\) denotes second quantization, not the nonadic holonomy. For the four-component spinorial module, the matrices are
\[
 \begin{array}{c|cccc}
 abc&012&013&023&123\\ \hline
 M_{abc}=A_Dg^ag^bg^c&
 iJ\otimes R&iJ\otimes Z&iI_2\otimes J&iZ\otimes J
 \end{array}.
\]
Multiplying \eqref{hmtcuatro:ant:clifford} gives these matrices; they are Hermitian, traceless and square to \(I_4\).

\begin{proposition}[Normal contraction and preservation of cross-terms]
\label{hmtcuatro:ant:orden-normal} For every matrix \(M\),
\[
 :d\Gamma(M)^2:
 =d\Gamma(M)^2-d\Gamma(M^2).
\]
In the four-component module,
\[
 \widehat{\mathscr Q}_{\rm sp}
 =\sum_{a=1}^4s_a d\Gamma(M_a)^2+2\widehat N,
 \qquad s=(-1,-1,-1,+1),\quad \widehat N=d\Gamma(I_4).
\]
The operator is self-adjoint and preserves occupation. In a sum of multiplets, the same first identity is used with the total matrices \(\widetilde M_a\); the contractions between species are retained.
\end{proposition}
\begin{proof}
In \(d\Gamma(M)^2\), the relation \(a_sa_t^\dagger=\delta_{st}-a_t^\dagger a_s\) separates the contraction \(d\Gamma(M^2)\) from the term with two creation and two annihilation operators. The latter is, by definition, the normal product. Since \(M_a^2=I_4\) and \(\sum_a s_a=-2\), the contraction subtraction equals \(+2\widehat N\). Hermiticity of the matrices makes the summands self-adjoint, and each monomial preserves occupation.

If \(\widetilde M=\bigoplus_s M_s\), then \(d\Gamma(\widetilde M)=\sum_s d\Gamma(M_s)\) and \(d\Gamma(\widetilde M^2)=\sum_s d\Gamma(M_s^2)\). Bilinears of distinct multiplets commute; squaring the first sum leaves \(2d\Gamma(M_s)d\Gamma(M_t)\) for \(s<t\). Subtracting an occupation does not remove these cross-terms. This is the operator version of \eqref{hmtcuatro:ant:cruces-especies}.
\end{proof}

Invariance under internal unitary changes follows because they act on the material indices and intertwine the current matrices: second quantization transports the full formula by conjugation. The density and its evaluation on a state also use the declared cell volume and material state. The mean current and its second moment are not identified. These data complete the background used by the joint action, without introducing a second torsion elimination or modifying its boundary conditions.
